% Avsnitt 2.9, ur lectures/L02/appendix/b_exercises.md.

\section{Övningar}
\label{c2:sec:exercises}

\begin{aside}
\textbf{Så kontrollerar du ditt arbete.} Varje övning nedan som ber dig skriva en VHDL-modul har
en utdelad självkontrollerande testbänk under \file{lectures/L02/exercises/}. Skriv din modul i
dess katalog \file{exercises/<module>/}, med det entitetsnamn och den \textbf{portordning}
övningen anger, och kör den sedan med GHDL \textemdash{} se \secref{c2:sec:testbenches} för de tre
kommandona.

Inget FPGA-kort behövs till någon övning. Stegen med Quartus-syntes och kortprogrammering
demonstreras under passet; ditt jobb efteråt är att få VHDL-koden rätt, och testbänken är hur du
bekräftar det.
\end{aside}

% ------------------------------------------------------------------------------------------
\exerciseset{Karnaughdiagram}

\exercise{Tre variabler}{Krets}
\label{c2:ex:kmap3}
Härled en minimerad ekvation för \code{X} ur Karnaughdiagrammet nedan, realisera sedan motsvarande
grindnät och simulera det i CircuitVerse:

\begin{narrowtable}{cc}
  \thead{ABC} & \thead{X} \\
  \midrule
  000 & 1 \\
  001 & 1 \\
  010 & 0 \\
  011 & 0 \\
  100 & 1 \\
  101 & 1 \\
  110 & 0 \\
  111 & 0 \\
\end{narrowtable}

\note[Tips]Lägg \code{AB} nedåt i raderna i Graykodsordning (\code{00, 01, 11, 10}) och \code{C} i
sidled i kolumnerna, precis som i det genomarbetade exemplet i \secref{c2:sec:kmap}, innan du
börjar gruppera \code{1}:or.

\exercise{Fyra variabler}{Krets}
\label{c2:ex:kmap4}
Härled en minimerad ekvation för \code{X} ur Karnaughdiagrammet nedan, realisera sedan motsvarande
grindnät och simulera det i CircuitVerse:

\begin{narrowtable}{cc}
  \thead{ABCD} & \thead{X} \\
  \midrule
  0000 & 0 \\
  0001 & 1 \\
  0010 & 0 \\
  0011 & 1 \\
  0100 & 0 \\
  0101 & 0 \\
  0110 & 0 \\
  0111 & 0 \\
  1000 & 0 \\
  1001 & 1 \\
  1010 & 0 \\
  1011 & 1 \\
  1100 & 1 \\
  1101 & 0 \\
  1110 & 1 \\
  1111 & 0 \\
\end{narrowtable}

\note[Tips]Med fyra variabler lägger du två av dem (i Graykod) på varje axel, så att du får ett
4x4-rutnät. Grupper kan fortfarande gå runt varje kant av rutnätet, inte bara vänster\--höger; kom
ihåg att kontrollera överkant och underkant också.

% ------------------------------------------------------------------------------------------
\exerciseset{Kretsar i VHDL}

\exercise{\code{xyz_logic}: tre utgångar ur fyra ingångar}{VHDL}
\label{c2:ex:xyz}
Ett grindnät har fyra ingångar, \code{ABCD}, och tre utgångar, \code{XYZ}. Det ges av
sanningstabellen nedan:

\begin{narrowtable}{cc}
  \thead{ABCD} & \thead{XYZ} \\
  \midrule
  0000 & 001 \\
  0001 & 010 \\
  0010 & 001 \\
  0011 & 011 \\
  0100 & 100 \\
  0101 & 110 \\
  0110 & 100 \\
  0111 & 111 \\
  1000 & 101 \\
  1001 & 110 \\
  1010 & 101 \\
  1011 & 111 \\
  1100 & 000 \\
  1101 & 010 \\
  1110 & 000 \\
  1111 & 011 \\
\end{narrowtable}

\xpart{a} Behandla \code{X}, \code{Y} och \code{Z} som tre separata funktioner av \code{ABCD} med
en utgång var, och härled en minimerad ekvation för var och en via ett Karnaughdiagram.

\xpart{b} Realisera det resulterande grindnätet och simulera det i CircuitVerse; stäm av alla tre
utgångarna mot sanningstabellen ovan.

\xpart{c} Implementera konstruktionen i VHDL, i en egen modul med fyra \code{std_logic}-ingångar
och tre \code{std_logic}-utgångar. \Secref{c1:sec:vhdl} och \secref{c2:sec:tovhdl} visar mönstret
för att göra konkurrenta VHDL-tilldelningar av en härledd ekvation.

\xpart{d} Kör testbänken och bekräfta att alla 16 raderna går igenom.

Om en rad fallerar namnger assertion-meddelandet ingångskombinationen:
\begin{itemize}
  \item Gå tillbaka till Karnaughdiagrammet för den av \code{X}, \code{Y}, \code{Z} som är fel.
  \item En enda fel utgång på en enda rad är nästan alltid en felgrupperad \code{1} i ett diagram.
\end{itemize}

\note[Tips]Två av de tre utgångarna här visar sig bero på bara en eller två av de fyra ingångarna
när de väl minimerats. Anta inte att varje utgång behöver alla fyra.

\modulefig[0.52]{lectures/L02/appendix/images/xyz_logic.png}

\begin{selfcheck}
\textbf{Självkontroll.} Döp din VHDL-entitet till \code{xyz_logic}, med portarna \code{a},
\code{b}, \code{c}, \code{d} (in) och \code{x}, \code{y}, \code{z} (out), deklarerade i den
ordningen; dess testbänk ligger i \file{exercises/xyz_logic/} och kontrollerar alla 16 raderna i
sanningstabellen ovan.
\end{selfcheck}

% ------------------------------------------------------------------------------------------
\exerciseset{Multiplexrar}

\exercise{4-till-1-muxen för hand}{Krets}
\label{c2:ex:mux4hand}
En 4-till-1-multiplexer har fyra dataingångar, \code{A}--\code{D}, två väljarbitar, \code{S[1:0]},
och en utgång, \code{X}.

Låt väljaren namnge sin dataingång direkt: \code{S[1:0] = "00"} väljer \code{A}, \code{"01"} väljer
\code{B}, \code{"10"} väljer \code{C} och \code{"11"} väljer \code{D}. Det är den avbildning
Övning~\ref{c2:ex:mux4} implementerar.

\xpart{a} Skriv ut multiplexerns sanningstabell (fyra rader, en per väljarkombination).

\xpart{b} Härled summa-av-produkter-ekvationen för \code{X} uttryckt i \code{A}--\code{D} och
\code{S[1:0]}. \Secref{c2:sec:mux} gör likadant för 8-till-1-fallet; halvera det här.

\xpart{c} Bygg och simulera nätet i CircuitVerse, och bekräfta att \code{X} följer varje
dataingång i tur och ordning när du ändrar \code{S[1:0]}.

Spara ditt svar: Övning~\ref{c2:ex:mux4} nedan implementerar precis den här kretsen i VHDL, när
\secref{c2:sec:process} väl har gett dig de \code{process}- och \code{case}-satser den behöver.

\exercise{\code{mux2} med \code{if}/\code{else}}{VHDL}
\label{c2:ex:mux2}
\Secref{c2:sec:process} skriver 2:1-multiplexern med en \code{case}-sats. Skriv samma krets en
gång till med en \code{if}/\code{else} i stället, och övertyga dig själv om att de två beskriver
identisk hårdvara.

Skriv en entitet med namnet \code{mux2} med:

\begin{booktable}{@{}p{3.6em}p{3.4em}p{6.4em}L@{}}
  \thead{Port} & \thead{Riktn.} & \thead{Typ} & \thead{Beskrivning} \\
  \midrule
  \code{d0} & in & \code{std_logic} & Dataingång. \\
  \code{d1} & in & \code{std_logic} & Dataingång. \\
  \code{sel} & in & \code{std_logic} & Väljare. \\
  \code{x} & out & \code{std_logic} & Den valda dataingången. \\
\end{booktable}

Väljaren kopplar fram den dataingång vars nummer den namnger:
\begin{itemize}
  \item \code{sel = '0'} väljer \code{d0}.
  \item \code{sel = '1'} väljer \code{d1}.
  \item Observera att det är omvänt mot \code{case}-exemplet i \secref{c2:sec:process}, som väljer
    sin \emph{andra} ingång \code{b} vid \code{'0'}. Vilken ingång ett väljarvärde namnger är en
    konvention, inte en regel, så läs av det ur specifikationen varje gång i stället för att anta.
\end{itemize}

Implementera multiplexern med en processbaserad arkitektur:
\begin{itemize}
  \item Ta med \code{d0}, \code{d1} och \code{sel} i känslighetslistan.
  \item Använd en \code{if}/\code{else}-sats på \code{sel}.
  \item Tilldela \code{x} i varje gren.
  \item Använd inte:
  \begin{itemize}
    \item En villkorlig konkurrent tilldelning (\code{x <= d1 when sel = '1' else d0;}).
    \item En \code{with...select}-sats: en konkurrent form av \code{case} som boken inte går
      igenom, nämnd bara för att du inte ska sträcka dig efter den om du mött den någon annanstans.
  \end{itemize}
\end{itemize}

\modulefig[0.5]{lectures/L02/appendix/images/mux2.png}

\begin{selfcheck}
\textbf{Självkontroll.} Döp din entitet till \code{mux2}, med portarna \code{d0}, \code{d1},
\code{sel} (in) och \code{x} (out), deklarerade i den ordningen; dess testbänk ligger i
\file{exercises/mux2/} och kontrollerar alla åtta kombinationer av de tre ingångarna.
\end{selfcheck}

\exercise{\code{mux4}}{VHDL}
\label{c2:ex:mux4}
Utöka Övning~\ref{c2:ex:mux2} till en 4:1-multiplexer med namnet \code{mux4}.

Det här är kretsen du härledde, ritade och simulerade för hand i Övning~\ref{c2:ex:mux4hand}.
Härled \code{case}-grenarna på nytt själv för fyra ingångar och två väljarbitar i stället för att
kopiera 8-till-1-exemplet från \secref{c2:sec:process}.

Entiteten måste ha:

\begin{booktable}{@{}p{3.6em}p{3.4em}p{10.8em}L@{}}
  \thead{Port} & \thead{Riktn.} & \thead{Typ} & \thead{Beskrivning} \\
  \midrule
  \code{d0} & in & \code{std_logic} & Dataingång. \\
  \code{d1} & in & \code{std_logic} & Dataingång. \\
  \code{d2} & in & \code{std_logic} & Dataingång. \\
  \code{d3} & in & \code{std_logic} & Dataingång. \\
  \code{sel} & in & \code{std_logic_vector(1 downto 0)} & Väljare. \\
  \code{x} & out & \code{std_logic} & Den valda dataingången. \\
\end{booktable}

Som i Övning~\ref{c2:ex:mux2} namnger väljaren sin dataingång direkt: \code{"00"} väljer \code{d0},
\code{"01"} väljer \code{d1}, \code{"10"} väljer \code{d2} och \code{"11"} väljer \code{d3}.
Återigen är det den omvända avbildningen mot \secref{c2:sec:process}:s genomarbetade
8-till-1-exempel, där \code{"000"} kopplar fram \code{inputs(7)}. Härled dina \code{case}-grenar ur
de fyra raderna ovan, inte ur den listningen.

Implementera multiplexern med en process:
\begin{itemize}
  \item Ta med alla dataingångar och \code{sel} i känslighetslistan.
  \item Använd en \code{case}-sats på \code{sel}, en gren per väljarvärde ovan.
  \item Ta med en \code{when others}-gren, och välj ett säkert utgångsvärde för väljaringångar som
    innehåller värden som \code{'X'}, \code{'U'} eller \code{'Z'}.
  \item Ersätt inte \code{case}-satsen med en \code{if}/\code{elsif}-kedja.
\end{itemize}

\modulefig[0.56]{lectures/L02/appendix/images/mux4.png}

\begin{selfcheck}
\textbf{Självkontroll.} Döp din entitet till \code{mux4}, med portarna \code{d0}, \code{d1},
\code{d2}, \code{d3}, \code{sel} (\code{std_logic_vector(1 downto 0)}) och \code{x} (out),
deklarerade i den ordningen; dess testbänk ligger i \file{exercises/mux4/} och sveper alla fyra
väljarvärdena mot alla sexton dataingångskombinationerna.
\end{selfcheck}

% ------------------------------------------------------------------------------------------
\exerciseset{Submoduler och 7-segmentsdisplayer}

Övning~\ref{c2:ex:display} och \ref{c2:ex:hexdisplay} bygger en konstruktion i två halvor: en
modul som driver en enda 7-segmentsdisplay, och en toppnivå som använder två kopior av den för att
visa en hexadecimal siffra på var och en. Gör dem i ordning: Övning~\ref{c2:ex:hexdisplay}
instansierar det du skriver i Övning~\ref{c2:ex:display}.

En 7-segmentsdisplay är sju lysdioder, namngivna \code{a} till \code{g}, arrangerade runt en åtta.
Tänd rätt delmängd så får du en siffra. Det här är utseendet, och det är det du behöver för att
kunna räkna ut koderna nedan:

\begin{textdrawing}
      aaaa
     f    b
     f    b
      gggg
     e    c
     e    c
      dddd
\end{textdrawing}

Så \code{1} tänder \code{b} och \code{c} (de två till höger), och \code{7} tänder \code{a},
\code{b} och \code{c}. Att jämföra koderna som ges för de två siffrorna är ett snabbt sätt att
bekräfta att du har utseendet rättvänt innan du börjar härleda \code{A} till \code{F}.

På DE0-CV är displayerna kopplade \textbf{aktivt låga}: att driva ett segment till \code{'0'}
tänder det, och \code{'1'} släcker det. Det är därför koden för \code{8}, som tänder varje
segment, är \code{"0000000"}, och varför en släckt display är \code{"1111111"}.

Genomgående avbildas \code{hex(6 downto 0)} på segmenten \code{g}, \code{f}, \code{e}, \code{d},
\code{c}, \code{b}, \code{a}: bit 6 är \code{g} och bit 0 är \code{a}.

\exercise{\code{display}}{VHDL}
\label{c2:ex:display}
Skriv en modul med namnet \code{display} som driver en 7-segmentsdisplay med en hexadecimal
siffra.

Entiteten måste ha:

\begin{booktable}{@{}p{5em}p{3.4em}p{10.8em}L@{}}
  \thead{Port} & \thead{Riktn.} & \thead{Typ} & \thead{Beskrivning} \\
  \midrule
  \code{number} & in & \code{std_logic_vector(3 downto 0)} & Värdet som ska visas, \code{0000} till
    och med \code{1111}. \\
  \code{hex} & out & \code{std_logic_vector(6 downto 0)} & Segmentkoden för det värdet. \\
\end{booktable}

Att mata \code{number} värdet \code{0111} måste alltså lägga koden för \code{7} på \code{hex}, och
att mata den \code{1110} måste lägga koden för \code{E} på \code{hex}.

Implementera den med en process:
\begin{itemize}
  \item Lägg \code{number} i känslighetslistan.
  \item Använd en \code{case}-sats på \code{number}, en gren per värde.
  \item Ta med en \code{when others}-gren, och släck displayen i den.
\end{itemize}

Du får koderna för \code{0} till \code{9}. Räkna ut \code{A} till \code{F} själv, utifrån
segmentlayouten ovan:

\begin{vhdlcode}
constant DISPLAY_0  : std_logic_vector(6 downto 0) := "1000000";
constant DISPLAY_1  : std_logic_vector(6 downto 0) := "1111001";
constant DISPLAY_2  : std_logic_vector(6 downto 0) := "0100100";
constant DISPLAY_3  : std_logic_vector(6 downto 0) := "0110000";
constant DISPLAY_4  : std_logic_vector(6 downto 0) := "0011001";
constant DISPLAY_5  : std_logic_vector(6 downto 0) := "0010010";
constant DISPLAY_6  : std_logic_vector(6 downto 0) := "0000010";
constant DISPLAY_7  : std_logic_vector(6 downto 0) := "1111000";
constant DISPLAY_8  : std_logic_vector(6 downto 0) := "0000000";
constant DISPLAY_9  : std_logic_vector(6 downto 0) := "0010000";
-- Todo: add DISPLAY_A through DISPLAY_F here.
constant DISPLAY_OFF: std_logic_vector(6 downto 0) := "1111111";
\end{vhdlcode}

Två av de sex är gemener på en 7-segmentsdisplay, eftersom deras versalformer skulle vara omöjliga
att skilja från en siffra: \code{b} skulle läsas som \code{8} och \code{d} skulle läsas som
\code{0}. \code{A}, \code{C}, \code{E} och \code{F} är versaler.

Alla sexton koderna sitter också i \file{display_tb.vhd}, eftersom en självkontrollerande testbänk
inte kan kontrollera en utgång utan att veta vad den borde vara. Så det här är inget pussel du kan
bli utelåst från. Härled de sex ändå: det tar ett par minuter med segmentdiagrammet, och att läsa
av dem ur testbänken lär dig ingenting som \code{case}-satsen inte gör. Jobbet i den här övningen
är \code{case}-satsen med sexton grenar, inte uppslagstabellen.

\modulefig[0.68]{lectures/L02/appendix/images/display.png}

\begin{selfcheck}
\textbf{Självkontroll.} Döp din entitet till \code{display}, med portarna \code{number}
(\code{std_logic_vector(3 downto 0)}) och \code{hex} (out, \code{std_logic_vector(6 downto 0)}),
deklarerade i den ordningen; dess testbänk ligger i \file{exercises/display/} och kontrollerar alla
sexton värdena mot hela kodtabellen, och namnger siffran och båda koderna när en är fel.
\end{selfcheck}

\exercise{\code{hex_display}}{VHDL}
\label{c2:ex:hexdisplay}
Skriv en modul med namnet \code{hex_display} som visar ett tvåsiffrigt hexadecimalt tal på två
displayer, med två instanser av din \code{display} från Övning~\ref{c2:ex:display}.

Entiteten måste ha:

\begin{booktable}{@{}p{4.6em}p{3.4em}p{10.8em}L@{}}
  \thead{Port} & \thead{Riktn.} & \thead{Typ} & \thead{Beskrivning} \\
  \midrule
  \code{input} & in & \code{std_logic_vector(7 downto 0)} & Åtta strömbrytare: de fyra översta är
    den vänstra siffran, de fyra nedersta den högra. \\
  \code{hex1} & out & \code{std_logic_vector(6 downto 0)} & Segmentkod för den vänstra siffran. \\
  \code{hex0} & out & \code{std_logic_vector(6 downto 0)} & Segmentkod för den högra siffran. \\
\end{booktable}

När den fungerar dyker talet på \code{input(7 downto 4)} upp på \code{hex1}, och talet på
\code{input(3 downto 0)} på \code{hex0}.

Bygg den med submoduler, enligt \secref{c2:sec:submodules}:
\begin{itemize}
  \item Skapa två instanser av \code{display}, med etiketterna \code{display1} och \code{display0}.
  \item Koppla \code{input(7 downto 4)} och \code{hex1} till \code{display1}.
  \item Koppla \code{input(3 downto 0)} och \code{hex0} till \code{display0}.
  \item Använd positionell \code{port map}, som överallt annars i den här boken.
  \item Skriv ingen egen logik i \code{hex_display}. Om du kommer på dig själv med att skriva en
    \code{case}-sats här duplicerar du Övning~\ref{c2:ex:display} i stället för att återanvända
    den.
\end{itemize}

Kopiera in din \file{display.vhd} från Övning~\ref{c2:ex:display} i övningskatalogen, eftersom
\code{hex_display} inte kan analyseras utan den. Det här är bokens första konstruktion byggd av
mer än en fil, så namnge \file{display.vhd} först, före den fil som instansierar den:

\begin{shell}
ghdl -a --std=93 display.vhd hex_display.vhd hex_display_tb.vhd
ghdl -e --std=93 hex_display_tb
ghdl -r --std=93 hex_display_tb --assert-level=error --stop-time=10ms
\end{shell}

\Secref{c2:sub:multifile} förklarar varför ordningen spelar roll.

\modulefig[0.68]{lectures/L02/appendix/images/hex_display.png}

\begin{selfcheck}
\textbf{Självkontroll.} Döp din entitet till \code{hex_display}, med portarna \code{input}
(\code{std_logic_vector(7 downto 0)}), \code{hex1} (out) och \code{hex0} (out, båda
\code{std_logic_vector(6 downto 0)}), deklarerade i den ordningen; dess testbänk ligger i
\file{exercises/hex_display/} och sveper alla 256 ingångsvärdena, och kontrollerar varje halva
oberoende så att en konstruktion som kastar om de två siffrorna fångas i stället för att gå igenom
på de värden där båda halvorna råkar stämma överens.
\end{selfcheck}

\exercise{En siffra som grindnät}{Krets}
\label{c2:ex:segments}
\emph{Om du har tid.} Realisera en siffra som ett grindnät, så som \chapref{c1:ch} och
\secref{c2:sec:kmap} lät dig göra:
\begin{itemize}
  \item Ta de fyra ingångarna och ett segment, säg \code{a}, och härled den booleska ekvationen
    för det ur sanningstabellen med 16 rader. Ett Karnaughdiagram är värt att använda här: de
    flesta segment minimeras rejält.
  \item Upprepa för så många av de sex återstående segmenten som du har tålamod med, och bygg
    resultatet i CircuitVerse.
  \item Kopiera sedan hela grindnätet för att driva en andra display.
\end{itemize}

Den kopian är poängen med övningen. Att duplicera ett block grindar är precis vad det gjorde att
instansiera \code{display} en andra gång i Övning~\ref{c2:ex:hexdisplay}, och det är värt att ha
gjort en gång för hand för att se vad den enda raden VHDL står för.

% ------------------------------------------------------------------------------------------
\exerciseset{Att namnge ett mellanresultat}

\exercise{\code{combo_logic}, skriven två gånger}{VHDL}
\label{c2:ex:combo}
Skriv en entitet med namnet \code{combo_logic} som implementerar den booleska funktionen

\[ x = ab + c'd \]

\Secref{c2:sec:tovhdl} hävdade att det inte kostar något och ger läsbarhet att namnge ett
deluttryck med en intern signal. Den här övningen är där du bekräftar det snarare än tar det på
tro: du skriver funktionen två gånger, en gång som ett enda uttryck och en gång med termen $c'd$
utbruten i en intern signal med namnet \code{y}.

En signal deklarerad inuti en arkitektur är en ledning. Den är ingen variabel, den lagrar
ingenting, och att namnge ett deluttryck med en sådan lägger varken till en grind, en fördröjning
eller en vippa i kretsen. Den ger bara en del av nätet ett namn, så att du, när du läser det
senare, kan se vad den delen är till för. Båda arkitekturerna nedan syntetiseras till samma
grindar, och den utdelade testbänken går igenom mot vilken som helst av dem.

Entiteten måste ha portarna \code{a}, \code{b}, \code{c}, \code{d} (in) och \code{x} (out), alla
\code{std_logic}.

\xpart{a} Implementera hela funktionen med en enda konkurrent signaltilldelning: deklarera inga
interna signaler, och översätt varje boolesk operation till motsvarande VHDL-operator.

\xpart{b} Skriv om arkitekturen med en intern signal med namnet \code{y}, som representerar termen
$c'd$, och tilldela \code{x} det slutliga uttrycket. Det här är samma form som
\code{gate_network} i \secref{c2:sec:tovhdl}, som namnger sin \code{AND}-term \code{ab}.

Bekräfta att de två arkitekturerna beskriver samma krets.

\note[Tips]Två beskrivningar representerar samma kombinatoriska krets när de producerar samma
utgång för varje möjlig ingångskombination. Med fyra ingångar finns 16 möjliga
ingångskombinationer: skriv den fullständiga sanningstabellen och verifiera att båda arkitekturerna
producerar samma värde på \code{x} på varje rad.

\modulefig[0.52]{lectures/L02/appendix/images/combo_logic.png}

\begin{selfcheck}
\textbf{Självkontroll.} Döp din entitet till \code{combo_logic}, med portarna \code{a}, \code{b},
\code{c}, \code{d} (in) och \code{x} (out), deklarerade i den ordningen; dess testbänk ligger i
\file{exercises/combo_logic/} och kontrollerar alla 16 raderna. Den går igenom mot vilken som helst
av arkitekturerna, vilket är poängen med jämförelsen ovan.
\end{selfcheck}
